% GENERATED FILE. Edit canonical lessons, then run npm run generate:books.
% canonical-source-hash: fnv1a64:a4a984a103b02ad8
% canonical-lessons: HI-C121-dhyanrakhie, HI-C121-sambhalkar, HI-C121-shubhyatra, HI-C121-phiraiega, HI-C121-chalie

\chapter{Take Care}
\label{ch:hi-take-care}
\hlchaptermodality{121}

\begin{chapteropening}
\textbf{By the end of this chapter:} \emph{I can tell someone to take care and go carefully, wish them a good journey, ask them to come again, and say let's go.}

Complete HI-C121-chalie: leave a visit well.
\end{chapteropening}

\section[\texorpdfstring{\hi{ध्यान} \hi{रखिए} (dhyān rakhie)}{ध्यान रखिए (dhyān rakhie)}]{\hi{ध्यान} \hi{रखिए} (dhyān rakhie) — take care}

\label{lesson:HI-C121-dhyanrakhie}



\begin{quote}
Before the new one: say the Hindi for a person, then the Hindi for nice to meet you.
\end{quote}

\subsection*{You'll want to know: \hi{ध्यान} \hi{रखिए}}
\textbf{\hi{ध्यान} \hi{रखिए}} (\emph{dhyān rakhie}) — \textquotedblleft{}take care\textquotedblright{}, literally \textquotedblleft{}keep attention\textquotedblright{}.

\textbf{\hi{अपना} \hi{ध्यान} \hi{रखिए}} — \textquotedblleft{}take care of yourself\textquotedblright{}, the kindest goodbye.

\begin{cousinweb}
\textbf{\hi{ध्यान} \hi{रखिए}} (\emph{dhyān rakhie}) is \textquotedblleft{}keep attention\textquotedblright{}. \emph{dhyān} is Sanskrit \emph{dhyāna}, \textquotedblleft{}meditation\textquotedblright{}, a tatsama (CDIAL 6889), and \emph{rakhnā}, \textquotedblleft{}to keep\textquotedblright{}, comes from \emph{rakṣati}, \textquotedblleft{}guards\textquotedblright{} (CDIAL 10547).
\end{cousinweb}

\begin{grammarlens}[title={What you've built}]
The kindest thing to say at the door.
\end{grammarlens}

\subsection*{Your turn}
\begin{itemize}
\raggedright
  \item \emph{Say it:} \emph{dhyān rakhie}
  \item \emph{Say it:} \emph{dhyān rakhie}, once more
  \item \emph{Say it:} say \emph{apnā dhyān rakhie}
  \item \emph{Recall:} say the Hindi for a person, then the Hindi for nice to meet you, then say \emph{dhyān rakhie} again
\end{itemize}

\begin{tcolorbox}[breakable,colback=teal!4,colframe=teal!35!black,title={Before you move on}]
What does \hi{ध्यान} \hi{रखिए} mean? (\textquotedblleft{}Take care\textquotedblright{}.) Say it once more.
\end{tcolorbox}
\section[\texorpdfstring{\hi{संभलकर} (sambhalkar)}{संभलकर (sambhalkar)}]{\hi{संभलकर} (sambhalkar) — carefully}

\label{lesson:HI-C121-sambhalkar}



\begin{quote}
Before the new one: say the Hindi for nice to meet you, then the Hindi for take care.
\end{quote}

\subsection*{You'll want to know: \hi{संभलकर}}
\textbf{\hi{संभलकर}} (\emph{sambhalkar}) — \textquotedblleft{}carefully\textquotedblright{}, \textquotedblleft{}steadily\textquotedblright{}.

\textbf{\hi{संभलकर} \hi{जाना}} — \textquotedblleft{}go carefully\textquotedblright{}, said to someone setting out after dark or in the rain.

\begin{cousinweb}
\textbf{\hi{संभलकर}} (\emph{sambhalkar}) is \emph{sambhalnā}, \textquotedblleft{}to steady oneself\textquotedblright{}, with \emph{-kar}, \textquotedblleft{}having\textquotedblright{}. CDIAL lists the verb under Sanskrit \emph{saṁbharati}, \textquotedblleft{}brings together\textquotedblright{}, and its causing form \textquotedblleft{}maintains\textquotedblright{} (CDIAL 12961).
\end{cousinweb}

\begin{grammarlens}[title={What you've built}]
A second way to wish someone safe.
\end{grammarlens}

\subsection*{Your turn}
\begin{itemize}
\raggedright
  \item \emph{Say it:} \emph{sambhalkar}
  \item \emph{Say it:} \emph{sambhalkar}, once more
  \item \emph{Say it:} say \emph{sambhalkar jānā}
  \item \emph{Recall:} say the Hindi for nice to meet you, then the Hindi for take care, then say \emph{sambhalkar} again
\end{itemize}

\begin{tcolorbox}[breakable,colback=teal!4,colframe=teal!35!black,title={Before you move on}]
What does \hi{संभलकर} mean? (\textquotedblleft{}Carefully\textquotedblright{}.) Say it once more.
\end{tcolorbox}
\section[\texorpdfstring{\hi{शुभ} \hi{यात्रा} (śubh yātrā)}{शुभ यात्रा (śubh yātrā)}]{\hi{शुभ} \hi{यात्रा} (śubh yātrā) — have a good journey}

\label{lesson:HI-C121-shubhyatra}



\begin{quote}
Before the new one: say the Hindi for take care, then the Hindi for carefully.
\end{quote}

\subsection*{You'll want to know: \hi{शुभ} \hi{यात्रा}}
\textbf{\hi{शुभ} \hi{यात्रा}} (\emph{śubh yātrā}) — \textquotedblleft{}have a good journey\textquotedblright{}.

\textbf{\hi{शुभ}} is the auspicious of \textbf{\hi{शुभ} \hi{रात्रि}}; \textbf{\hi{यात्रा}} is \textquotedblleft{}a journey\textquotedblright{}, a formal cousin of \textbf{\hi{सफ़र}}.

\begin{cousinweb}
\textbf{\hi{शुभ} \hi{यात्रा}} (\emph{śubh yātrā}) is two Sanskrit words taken over whole, tatsamas: \emph{śubha}, \textquotedblleft{}auspicious, good\textquotedblright{}, and \emph{yātrā}, \textquotedblleft{}journey\textquotedblright{}.
\end{cousinweb}

\begin{grammarlens}[title={What you've built}]
A wish for someone setting off.
\end{grammarlens}

\subsection*{Your turn}
\begin{itemize}
\raggedright
  \item \emph{Say it:} \emph{śubh yātrā}
  \item \emph{Say it:} \emph{śubh yātrā}, once more
  \item \emph{Say it:} say \emph{śubh yātrā!}
  \item \emph{Recall:} say the Hindi for take care, then the Hindi for carefully, then say \emph{śubh yātrā} again
\end{itemize}

\begin{tcolorbox}[breakable,colback=teal!4,colframe=teal!35!black,title={Before you move on}]
What does \hi{शुभ} \hi{यात्रा} mean? (\textquotedblleft{}Have a good journey\textquotedblright{}.) Say it once more.
\end{tcolorbox}
\section[\texorpdfstring{\hi{फिर} \hi{आइएगा} (phir āiegā)}{फिर आइएगा (phir āiegā)}]{\hi{फिर} \hi{आइएगा} (phir āiegā) — do come again}

\label{lesson:HI-C121-phiraiega}



\begin{quote}
Before the new one: say the Hindi for carefully, then the Hindi for a good journey.
\end{quote}

\subsection*{You'll want to know: \hi{फिर} \hi{आइएगा}}
\textbf{\hi{फिर} \hi{आइएगा}} (\emph{phir āiegā}) — \textquotedblleft{}do come again\textquotedblright{}.

\textbf{\hi{फिर}} is \textquotedblleft{}again\textquotedblright{}, and \textbf{\hi{आइएगा}} is an extra-polite \textbf{\hi{आइए}}, a request for later rather than now.

\begin{cousinweb}
\textbf{\hi{फिर} \hi{आइएगा}} (\emph{phir āiegā}) is \emph{phir}, \textquotedblleft{}again\textquotedblright{}, and \emph{āiegā}, the polite future of \textbf{\hi{आना}}, \textquotedblleft{}to come\textquotedblright{}. CDIAL lists \emph{phir} under a reconstructed \emph{*phirati}, \textquotedblleft{}turns\textquotedblright{}, whose old form became an adverb (CDIAL 9078).
\end{cousinweb}

\begin{grammarlens}[title={What you've built}]
An invitation to return.
\end{grammarlens}

\subsection*{Your turn}
\begin{itemize}
\raggedright
  \item \emph{Say it:} \emph{phir āiegā}
  \item \emph{Say it:} \emph{phir āiegā}, once more
  \item \emph{Say it:} say \emph{phir āiegā}
  \item \emph{Recall:} say the Hindi for carefully, then the Hindi for a good journey, then say \emph{phir āiegā} again
\end{itemize}

\begin{tcolorbox}[breakable,colback=teal!4,colframe=teal!35!black,title={Before you move on}]
What does \hi{फिर} \hi{आइएगा} mean? (\textquotedblleft{}Do come again\textquotedblright{}.) Say it once more.
\end{tcolorbox}
\section[\texorpdfstring{\hi{चलिए} (calie)}{चलिए (calie)}]{\hi{चलिए} (calie) — let's go, shall we}

\label{lesson:HI-C121-chalie}



\begin{quote}
Before the new one: say the Hindi for a good journey, then the Hindi for come again.
\end{quote}

\subsection*{You'll want to know: \hi{चलिए}}
\textbf{\hi{चलिए}} (\emph{calie}) — \textquotedblleft{}let's go\textquotedblright{}, \textquotedblleft{}shall we\textquotedblright{}.

It is \textbf{\hi{चलना}} with the respectful ending of \textbf{\hi{बोलिए}}, the verb inside \textbf{\hi{चलता} \hi{हूँ}}, \textquotedblleft{}I'll be off\textquotedblright{}.

\begin{cousinweb}
\textbf{\hi{चलिए}} (\emph{calie}) is the respectful command of \emph{calnā}, \textquotedblleft{}to move, go\textquotedblright{}, which comes from Sanskrit \emph{calati}, \textquotedblleft{}goes away\textquotedblright{} (CDIAL 4715).
\end{cousinweb}

\begin{grammarlens}[title={What you've built}]
Take care, carefully, a good journey, come again, let's go.
\end{grammarlens}

\subsection*{Your turn}
\begin{itemize}
\raggedright
  \item \emph{Say it:} \emph{calie}
  \item \emph{Say it:} \emph{calie}, once more
  \item \emph{Say it:} say \emph{calie}
  \item \emph{Recall:} say the Hindi for a good journey, then the Hindi for come again, then say \emph{calie} again
\end{itemize}

\begin{tcolorbox}[breakable,colback=teal!4,colframe=teal!35!black,title={Before you move on}]
What does \hi{चलिए} mean? (\textquotedblleft{}Let's go, shall we\textquotedblright{}.) Say it once more.
\end{tcolorbox}
